
\section{Procedural island generation}
\label{sec:coral-island-example-generation}

We propose a structured process for generating coral reef island terrains that takes the user's sketches and produces a complete 3D terrain model. This process begins with the creation of an initial height field based on user inputs, and applies wind deformation to introduce natural variations, and finally integrate coral reef features via subsidence and coral growth modelling.

The generation of coral reef islands with our system begins with two intuitive sketch-based inputs from the user: a top-view sketch and a profile-view sketch, which define the island's horizontal layout and vertical elevation profile. Users further refine the terrain by applying wind deformation strokes, approximating wind and waves effects on the island's shape. This combination of sketches and wind inputs gives users precise control over both the island's structure and its natural variations, such as irregular coastlines or concave features. In this section we demonstrate the usefulness of these sketches, describe the technical details, and present the method in \cref{alg:sketch-to-height-wind}.

\subsection{Initial height field generation}
\label{sec:coral-island-generation-initial}

\begin{figure}
    \autofitgraphics[width=0.90 \linewidth]{Cicia_island.png, Cicia_island-outlines.png}
    \autofitgraphics[width=0.90 \linewidth]{Cicia_island-3D.png}
    \caption[A real island and its different regions]{(Left) A real-world example of aerial image (and 3D visualisation on bottom) of an island (Cicia Island) may be segmented into regions. (Right) We represent the different regions boundaries.}
    \label{fig:coral-island-top-view-sketch}
\end{figure}

The top-view sketch defines the island's outline as seen from above. Using a simple drawing interface, the user can delineate the boundaries between key regions of the island, including the island itself, beaches, lagoons, and surrounding abysses (\cref{fig:coral-island-top-view-sketch}). Our system assumes that these regions are arranged concentrically around the centre of the island, with each boundary defined by a radial distance from the centre (\cref{fig:coral-island-parametric-distance}).

Each region's boundary is represented in polar coordinates $\left(\radius_\p, \angl_\p\right)$, with $\radius_\p$ indicating the radial distance from the island's centre and $\angl_\p$ representing the angular position. This polar representation allows us to map the user's sketch onto a circular framework, ensuring smooth transitions between regions and maintaining a coherent island layout.

In this sketch, the user defines the overall horizontal layout of the island, including the size and shape of each feature (\cref{fig:coral-island-procedural-height-only}). Variations in the outline are introduced by allowing the radial distances to vary with angle, ensuring that the island is not strictly symmetrical and introducing more natural, irregular shapes.

\begin{figure}[H]
    \autofitgraphics[]{binary-heights-input-only-outlines-2.png, binary-heights-output-only-heightmap-2.png, binary-heights-render-2.png}
    \autofitcaptions{Input sketch, Output height field, Result}
    \caption[Curve-based island generation using only outlines]{Using only the outlines of the island as an input sketch (left), we provide a height field depending on the region in which it relies (middle). The resulting terrain is displayed (right). }
    \label{fig:coral-island-procedural-height-only}
\end{figure}

\begin{figure}
    \autofitgraphics[]{profileFunction.pdf, schema_profile.jpg}
    \caption[Our profile function inspired by the typical cross-section representation of islands' reefs]{(Left) A profile function $\heightProfile$ is defined as a 1D function and represents the surface from the centre of the island to the abysses. (Right) The cross-section representation of an island is often represented as a 1D function defined using terrain features as landmarks.}
    \label{fig:coral-island-profile-function}
\end{figure}

The profile-view sketch defines the vertical elevation profile of the island along any radial direction, offering control over the island's height. In this view, the user specifies the elevation of different regions of the island, such as the island peak, beach, lagoon, abyss, and everything in between, by drawing the corresponding profile curve (\cref{fig:coral-island-profile-function}).

These regions correspond to key terrain transitions: the highest point of the island (centre), the island border, the beach, the lagoon, and the deep-sea abyss. We interpolate these milestones into a continuous 1D height function $\heightProfile(\distRegions_\p)$, where $\distRegions_\p$ represents a non-uniform region distance from the island's centre, and $h(\p) = \heightProfile(\distRegions_\p)$ gives the height at each point. This continuous profile ensures smooth elevation transitions across the island.

By combining the top-view and profile-view sketches, our method generates a coherent 3D terrain model that accurately reflects the user's design by revolution modelling.

The generation of the coral reef island terrain begins by transforming the user-defined top-view and profile-view sketches into a coherent 3D height field. This process combines the radial layout of the top-view sketch with the elevation information provided by the profile-view sketch, creating a terrain that accurately represents the desired features (i.e. the island, beaches, lagoons, and abyss).

For any point $\p \in \R^2$ on the terrain, we define its polar coordinates $(\radius_\p, \angl_\p)$ relative to the center of the island $\Center$ as:
\begin{align*}
    \radius_\p = \|\p - \Center\|, \qquad \angl_\p = \operatorname{atan2}(\p.y - \Center.y,\; \p.x - \Center.x).
\end{align*}
The radial distance $\radius_\p$ determines which region the point belongs to (island, beach, lagoon, reef, or abyss) based on the user-defined radial limits. Those outlines from the top-view sketch provide the exact boundaries between regions.

Each point's height is computed using the profile function $\heightProfile(\distRegions_\p)$. Instead of using the raw radial distance $\radius_\p$, we define a parametric region distance $\distRegions_\p$ that maps each point to a normalised position along the concentric regions (see \cref{fig:coral-island-parametric-distance}). The radial space is divided by user-defined boundaries $\Radius_0, \Radius_1, \dots, \Radius_n$, corresponding to the island centre, border, beach, lagoon, reef, and abyss.

\AltTextImageCancelled{
    When a point $\p$ lies between two boundaries $\Radius_{i}$ and $\Radius_{i+1}$, its parametric distance is
    \begin{align}
        \eqlabel{eq:coral-island-dist-regions}{Parametric distance}
        \distRegions_\p = i + \frac{\radius_\p - \Radius_{i}}{\Radius_{i + 1} - \Radius_{i}},
    \end{align}
    where $i$ is the index of the region containing $\p$ (i.e., $\Radius_i \le \radius_\p < \Radius_{i+1}$). This linear mapping stretches each region's radial span to the interval $[i, i+1[$, ensuring smooth interpolation across region boundaries.
    For any point $\p$, the height is finally computed as:
    \begin{align}
        \eqlabel{eq:coral-island-height-from-profile}{Height from profile funtion}
        h(\p) = \heightProfile(\distRegions_\p).
    \end{align}
}{outlines-top-view-x-bar.pdf, outlines-x-bar.pdf}{The $\tilde{x}_\p$ parameter is used to stretch the 1D height function $\heightProfile(x)$ to fit the distances from the centre to the outlines of each region defined in the top-view sketch.}{fig:coral-island-parametric-distance}{Stretching $\heightProfile(x)$ using $\tilde{x}_\p$}



This approach ensures that the height field accurately follows the elevation profile specified by the user while maintaining smooth transitions between different regions of the island.

The result is a height field that captures both the radial structure of the island (from the top-view sketch) and the vertical elevation profile (from the profile-view sketch), producing an organic representation of islands with smooth transitions between the key terrain features. Three slightly edited height functions on an identical island layout and their associated outputs are presented in \cref{fig:coral-island-procedural-smooth-heights}. To avoid perfectly smooth surfaces, we also apply a very low-amplitude Perlin noise perturbation to the final height field, adding fine-scale irregularities.

\begin{figure}
    \autofitgraphics[]{smooth-input-outline-heights-1.png, smooth-output-heights-1.png, smooth-render-1.png}
    \autofitgraphics[]{smooth-input-outline-heights-2.png, smooth-output-heights-2.png, smooth-render-2.png}
    \autofitgraphics[]{smooth-input-outline-heights-3.png, smooth-output-heights-3.png, smooth-render-3.png}
    \autofitcaptions{Input sketches, Output height field, Result}
    \caption[Curve-based generation of an island, varying the shape of the profile function]{Providing a smooth function between each region results in islands with plausible reliefs. We fixed the outlines (first column) while editing only the height function (second column) in order to produce, from top to bottom, a low island, a coral reef island, and finally an identical island without the reef.}
    \label{fig:coral-island-procedural-smooth-heights}
\end{figure}

\subsubsection{Wind and wave deformation}
\label{sec:coral-island-wind-deformation}

% In addition to the shape sketches, the user influences the shape of the island by defining a velocity field used to erode the initial terrain. This field simulates the effects of wind and wave erosion on the island's surface, introducing natural deformations such as coastline indentations, and more importantly allows the user to break the radial symmetry constraint.

In island environments, wind and waves reshape islands through coastal erosion, which removes and redistributes material over time. Simulating this physically requires sediment transport and hydrodynamics, which is both computationally intensive and difficult to control interactively. Instead, we approximate the morphological outcome of erosion with a deformation field: user-drawn strokes define a wind velocity field that warps the island, introducing large-scale asymmetries without explicitly removing matter. To account for the fact that not all regions respond equally to erosion, we also introduce a resistance function, which modulates deformation's strength across regions. For example, shallow coastal areas are strongly deformed, while the deep abyss or resistant volcanic core remain largely unaffected. This combination provides an efficient and intuitive proxy for the long-term shaping role of wind and waves.
Although real coastlines are shaped by both wind and wave processes, we represent them jointly through a single user-defined wind field, which serves as a controllable proxy for their combined morphological effect.

The wind field is represented as a series of strokes drawn by the user on a 2D canvas. Each stroke represents a parametric curve, where the direction and strength of the wind are encoded as a vector field. The user controls the wind's direction by drawing a set of curves, and we interpret them to create a velocity field that defines how the terrain is deformed.

As the user draws a stroke, a set of control points are generated along it to create a curve, with the option to adjust the stroke's width. The width of each stroke determines the area of influence around the curve, where wider strokes result in broader deformations of the terrain.
The deformation strength decreases with distance from the curve using a Gaussian falloff function using the stroke width as standard deviation. This ensure that the terrain transitions smoothly from deformed regions to non-deformed areas.

Once the strokes are applied to the velocity field, we displace each points of the terrain accordingly. The height field, originally generated from the user's sketches, is modified by the wind field to create non-radial features, breaking the initial radial symmetry and producing a more organic island shape (\cref{fig:coral-island-wind-stroke-edition}).


\begin{figure} 
    \autofitgraphics[]{
        1i-init-top-view.pdf, 
        1i-init-side-view.pdf, 
        1o_1.png, 
        % 1o_2.png, 
        1o_3-s.png}
    \autofitcaptions[]{Top view, Side view, Output height field, Result}
    \autofitgraphics[]{big-white.png, 3i-s.pdf, 3o_1.png, 3o_3-s.png}
    \autofitgraphics[]{big-white.png, 5i-s.pdf, 5o_1.png, 5o_3-s.png}
    \autofitgraphics[]{big-white.png, 7i-s.pdf, 7o_1.png, 7o_3-s.png}
    \autofitcaptions[]{~, Deformation, Output height field, Result}
    % \autofitgraphics[width=.6\linewidth]{3i-s.pdf, 3o_1.png, 3o_3-s.png}
    % \autofitgraphics[width=.6\linewidth]{5i-s.pdf, 5o_1.png, 5o_3-s.png}
    % \autofitgraphics[width=.6\linewidth]{7i-s.pdf, 7o_1.png, 7o_3-s.png}
    % \autofitcaptions[fontscale=1.5]{Deformation, Output height field, Result}
    % \autofitgraphics[width=.6\linewidth]{1i-init.pdf}
    % \autofitgraphics[trim={0 0 0 0.5cm}, clip]{
    %     1i-s.pdf, 
    %     % 2i-s.pdf, 
    %     3i-s.pdf, 
    %     % 4i-s.pdf, 
    %     5i-s.pdf, 
    %     % 6i-s.pdf, 
    %     7i-s.pdf, 
    %     % 8i-s.pdf, 
    %     % 9i-s.pdf
    % }
    % \autofitgraphics[]{
    %     1o_1.png,
    %     % 2o_1.png,
    %     3o_1.png,
    %     % 4o_1.png,
    %     5o_1.png,
    %     % 6o_1.png,
    %     7o_1.png,
    %     % 8o_1.png,
    %     % 9o_1.png,
    % }
    % \autofitgraphics[]{
    %     1o_3-s.png,
    %     % 2o_3-s.png,
    %     3o_3-s.png,
    %     % 4o_3-s.png,
    %     5o_3-s.png,
    %     % 6o_3-s.png,
    %     7o_3-s.png,
    %     % 8o_3-s.png,
    %     % 9o_3.png
    % }
    % % \autofitgraphics[]{
    % %     1o_2-s.png,
    % %     % 2o_2-s.png,
    % %     3o_2-s.png,
    % %     % 4o_2-s.png,
    % %     5o_2-s.png,
    % %     % 6o_2-s.png,
    % %     7o_2-s.png,
    % %     % 8o_2-s.png,
    % %     % 9o_2-s.png,
    % % }
    \caption[Island deformation through strokes]{Top row: A circular island layout and a given profile sketch creates a simple round island height field. Each following row shows a progression in an edition session, conserving the initial top-view and side-view sketch, by manipulating the deformation field (blue lines) with respectively 3, 7 and 20 user strokes (red arrows). The final result breaks the initial radial symmetry of the island layout. }
    \label{fig:coral-island-wind-stroke-edition}
\end{figure}

The deformation is controlled via a user-defined vector field representing the direction and strength of wind flows across the terrain. Users interact with our system by drawing strokes on a 2D canvas represented as centripetal Catmull-Rom splines $\curve$ \cite{Catmull1974}, a type of parametric curve that allows for smooth, continuous paths representing wind patterns. Each stroke defines a wind flow in the curve's direction $\curve'$, a strength $S_\curve$, and an effect width $\std$ (\cref{fig:coral-island-wind-from-strokes}); these wind flows are used to displace the terrain, approximating the gradual reshaping of the island due to wind and wave erosion. The strength of the displacement is controlled by a Gaussian scaling function $G_\std(x)$ that smoothly decreases with the distance from $\closestCp$, the closest point on the parametric curve $\curve$, as follows:

\begin{align}
    \eqlabel{eq:coral-island-deformation-from-stroke}{Displacement field from a single user stroke}
    \warp_\curve(\p) &= S_\curve \cdot \frac{\curve'(\closestCp)}{\norm{\curve'(\closestCp)} } \cdot G_\std\left(\norm{\p - \closestCp} \right), \\
    G_\std(x) &= \frac{1}{\std \sqrt{2\pi}} e^{-\frac{x^2}{2 \std^2}}. \nonumber
\end{align}

\begin{figure}
    \includegraphics[width = 0.8 \linewidth]{windByStrokes.pdf}
    \caption[Computation of a vector field from a single user wind stroke]{From the parametric curve defined by a user (red), we define the velocity field by considering the directional derivative $\curve'$ of the curve $\curve$ at the closest point $\closestCp$, modulated by a Gaussian distance function $G_{\std}(x)$ and scaled by the user-defined strength $S_\curve$.}
    \label{fig:coral-island-wind-from-strokes}
\end{figure}

The final deformation vector $\warp(\p)$ is computed as a sum of the influences from all strokes: 
\begin{align}
    \eqlabel{eq:coral-island-deformation-aggregation}{Displacement field from user strokes}
    \warp(\p) = \p + \sum_{\curve \in \text{curves}} \warp_\curve(\p) .
\end{align}

Once the deformation vector $\warp(\p)$ is computed, the terrain height at point $\p$ is adjusted by displacing $\p$ to a new point $\warp(\p)$.
We then compute the final height $h(\warp(\p)) = \heightProfile(\distRegions_{\p})$, or 
\begin{align}
    \eqlabel{eq:coral-island-height-deformation}{Height field with wind and wave deformations}
    \Tilde{h} = h \circ \warp^{-1}.
\end{align}

This process introduces variations in the terrain, distorting the coastline, creating concave regions, and breaking the original radial symmetry defined by the top-view and profile-view sketches.

\begin{figure}
    \includegraphics[width=0.45 \linewidth]{resistanceFunction.pdf}
    \caption[Resistance function]{The resistance function of the island is defined with a similar strategy as the $\heightProfile$ function. The resistance to erosion and deformation arises from multiple factors such as depth, materials, wind shadowing, biotic and abiotic factors, ... Modelling all these factors is complex. As such, using a user-defined approximation through a resistance function $\resistance$ allows for more control.}
    \label{fig:coral-island-resistance-function}
\end{figure}

To ensure that certain regions of the terrain such as deep-water areas remain relatively unaffected by the wind, a resistance function $\resistance(\distRegions_\p)$ is applied. The resistance function modulates the effect of deformation based on the previously computed piecewise parametric distance $\distRegions_\p$, with the same interaction means as the $\heightProfile$ function (\cref{fig:coral-island-resistance-function}).

% We define the resistance function $\resistance(\distRegions_\p)$ similarly to the profile function, and use it to modulate the magnitude of the displacement at each point. 
For example, regions near the coastline (such as the beach and lagoon) have lower resistance, allowing for significant deformation (emulating coastal erosion from wave energy), while regions farther away (such as the abyss) have higher resistance, limiting the wind and coastal erosion impact.

The deformation vector previously described is then scaled by the resistance function at each point $\p$, such that the final deformation vector becomes
\begin{align}
    \eqlabel{eq:coral-island-hindered-deformation}{Modulation of wind and wave field from resistance function}
    \Tilde{\warp}(\p) = \left(1 - \resistance(\distRegions_\p) \right) \cdot \warp(\p).
\end{align}

\begin{algorithm}[H]
    \caption{Sketches to height field with wind deformation.}
    \label{alg:sketch-to-height-wind}
    \DontPrintSemicolon
    \KwIn{Top-view boundaries $\{\Radius_0(\angl),\dots,\Radius_n(\angl)\}$, profile function $\heightProfile$, resistance $\resistance(\distRegions_\p)$, user strokes $\{\curve\}$ with width $\std$ and strength $S_\curve$}
    \KwOut{Deformed height field $\Tilde{h}$ and label map $\Tilde{I}$}

    \BlankLine

    \tcp{(1) Initial height field and label map from sketches}
    \ForEach{pixel $\p \in \R^2$}{
    $\radius_\p \gets \|\p - \Center\|$\;
    $\angl_\p \gets \operatorname{atan2}(\p.y-\Center.y,\;\p.x-\Center.x)$\;
    Find $i$ such that $\Radius_i(\angl_\p) \le \radius_\p < \Radius_{i+1}(\angl_\p)$\;
    $\distRegions_\p \gets i + \dfrac{\radius_\p - \Radius_i(\angl_\p)}{\Radius_{i+1}(\angl_\p) - \Radius_i(\angl_\p)}$\;
    $h(\p) \gets \heightProfile(\distRegions_\p)$\;
    $I(\p) \gets i$\;
    }

    \tcp{(2) Wind-field deformation from strokes}
    \ForEach{pixel $\p \in \R^2$}{
    $\warp(\p) \gets \mathbf{0}$\;
    \ForEach{stroke $\curve$}{
        $\closestCp \gets \arg\min_{\q \in \curve}\|\p-\q\|$ \tcp*{$\closestCp \gets$ closest point of the curve $\curve$ to the point $\p$}
        $\warp_\curve(\p) \gets S_\curve \cdot \dfrac{\curve'(\closestCp)}{\|\curve'(\closestCp)\|}\cdot G_\std\!\left(\|\p-\closestCp\|\right)$\tcp*{$G_\std(x)=\frac{1}{\std\sqrt{2\pi}}e^{-\frac{x^2}{2\std^2}}$}
        $\warp(\p) \gets \warp(\p) + \warp_\curve(\p)$\;
    }
    $\Tilde{\warp}(\p) \gets \left(1-\resistance(\distRegions_\p)\right)\cdot \warp(\p)$ \tcp*{Modulation from resistance function}
    }

    $\Tilde{h} \gets \Tilde{\warp}^{-1}\!\circ h$ \tcp*{$\Tilde{h}(\p) \gets h(\p - \Tilde{\warp})$ for backward warping}
    $\Tilde{I} \gets \Tilde{\warp}^{-1}\!\circ I$\;
    \Return $\Tilde{h}, \Tilde{I}$
\end{algorithm}

\begin{figure}
    \autofitgraphics[]{result_low_resistance.png, result_high_resistance.png}
    \autofitcaptions[]{Low resistance on reefs, High resistance on reefs}
    \caption[An island deformed with and without resistance]{(Left) Given a uniform velocity field and a resistance function similar to \cref{fig:coral-island-resistance-function}, the coasts are smoothly eroded while the interior of the island is almost unaffected. (Right) Modifying the resistance function to affect a strong resistance to borders emulates the effect of coast reinforcements.}
    \label{fig:coral-island-strong-low-resistance}
\end{figure}

The deformation process results in a modified height field where the terrain has been warped according to user-defined strokes; \cref{fig:coral-island-strong-low-resistance} illustrates the difference between non-resistant and resistant islands. This deformation introduces non-radial features, such as concave coastlines or irregularities along the beach and lagoon, making the island appear more natural and varied.

Both the height field and the label map (which tracks the terrain regions) are updated to reflect the deformation (\cref{fig:coral-island-wind-effect-result}). This ensures that the semantic information of the terrain remains consistent even after the terrain has been warped. The label map is deformed in the same way as the height field, preserving the logical structure of the island for further post-processing, such as texturing or mesh instancing.

In \cref{fig:coral-island-wind-stroke-edition}, a simple circular island is generated from the initial height field. By applying strokes along one side of the island, the deformation process can create concave regions along the coastline, making the shape more irregular and mimicking the effects of real-world wind and wave erosion. The resistance function ensures that while the beach and lagoon areas are deformed, the abyss remains largely unaffected as they are far from the wind and wave effective areas, preserving the island's overall structure.

\begin{figure}
    \setfitratios{1,1,1,1}
    \autofitcaptions[]{Top-view, Side-view, Deformation field, Resistance}
    \autofitgraphics[]{wind-deform-input-crop.png}
    % \autofitcaptions[]{}
    \setfitratios{1, 1}
    \autofitcaptions[]{Before deformation, After deformation}
    \autofitgraphics[]{wind-deform-output-original.png, wind-deform-labels-original.png, small-white.png, wind-deform-output.png, wind-deform-labels.png}
    % \setfitratios{1, 1, 1, 1, 1}
    \autofitcaptions[]{Height field, Label map, ~, Height field, Label map}
    \autofitgraphics[]{wind-deform-render-original.png, small-white.png,  wind-deform-render.png}
    % \setfitratios{1, 1}
    \autofitcaptions[]{Result, ~ , Result}
    \caption[Island deformation on specific regions using resistance function]{A top-view velocity field is defined from user-provided strokes (top, blue), in association with a resistance function (top, red); a height field is deformed accordingly. Left: original height field and render; right: altered results. The beach and lagoon regions are defined with low resistance, which is visible by having only these regions deformed in bottom results.}
    \label{fig:coral-island-wind-effect-result}
\end{figure}


\subsection{Coral reef modelling}
\label{sec:coral-island-coral-reef}

Once the terrain has been generated and deformed by the wind, we emulate the subsidence of the volcanic island and, in parallel, the growth of coral reefs. These processes reflect the long-term geological evolution of coral reef islands, where the volcanic island gradually sinks (subsides) while coral reefs grow upward to "keep up" with the sinking landmass.

\subsubsection{Subsidence}
\label{sec:coral-island-subsidence}

To account for the island subsidence due to tectonic activity we propose to scale the initial height field downward, modelling procedurally the effect of the volcanic island slowly sinking into the ocean. The user provides a subsidence rate $\subsidRate \in [0, 1]$, which represents the proportion by which the island has sunk. The subsidence is applied uniformly, meaning all terrain points on the island sink with the same intensity, regardless of their original height or location.

The subsided height field $\heightSubsid(\p)$ is computed by scaling the original height field $h(\p)$ with the subsidence factor:

\begin{align*}
    \heightSubsid(\p) = (1 - \subsidRate) \cdot h(\p).
\end{align*}

% This scaling reduces the overall height of the island, approximating how volcanic islands sink over time due to tectonic activity and erosion. The subsidence factor $\subsidRate$ is applied uniformly across the terrain, meaning that all points on the island experience the same degree of subsidence, regardless of their original height or location.

\subsubsection{Coral reef growth}
\label{sec:coral-island-reef-growth}

\begin{figure}
    \autofitgraphics[]{Reef_function_1.png, Reef_function_2.png, Reef_function_3.png}
    \caption[Coral reef growth from analytical model]{The modelling of the reef growth in our model is described by a piecewise function $\heightCoral$ which is flat in the lagoon, the crest and abyss, and follows a smoothstep function as transitions for the back reef and fore reef regions. The model is easily edited by tuning the height values and delimitations of the subregions. Here, three examples of reef growth function with different delimitations (vertical dotted lines), providing more or less smooth transitions between reef subregions.}
    \label{fig:coral-island-reef-function}
\end{figure}

As the volcanic island subsides, coral reefs grow upward to remain close to the water surface, following the "keep-up" strategy observed in most real-world coral formations. Coral growth is restricted to regions where the depth is within an optimal range for coral development, typically from the water surface to around 30 metres below before becoming much scarcer.

The coral reef features (reef crest, back reef, and fore reef) are modelled separately from the subsidence process. We generate a coral feature height field $\heightCoral(\p)$, which remains unaffected by the island's subsidence. %This height field ensures that coral regions remain near the water surface, even as the island sinks.

In our model, coral reef growth is entirely independent of the subsided terrain. Even as the volcanic island sinks, coral growth is driven only by the proximity of terrain to the water surface, ensuring that coral features always remain near the surface, irrespective of how much the island subsides.

We define distinct reef zones anchored at specific depths:

\begin{Itemize}
    \Item{} Reef crest near sea level: $h_{\text{crest}} = -2$\,m,
    \Item{} Back reef and lagoon: $h_{\text{back}} = -20$\,m,
    \Item{} Fore reef sloping to abyss: $h_{\text{abyss}} = -100$\,m.
\end{Itemize}

Each reef subregion is defined over a parametric domain $x \in [0, 1]$, with $x = 0$ the beginning of the reef region and $x = 1$ its end, directly inputting the parametric distance $x = \distRegions_\p - i_{\text{reef}}$ with $i_{\text{reef}}$ the ID of the reef region, resulting in a linear interpolation between the begining and the end of the reef region in the top-view sketch. For instance:
\begin{Itemize}
    \Item{} Back reef: $x_{\text{back,start}} = 0$, $x_{\text{back,end}} = 0.5$,
    \Item{} Reef crest: $x_{\text{crest,start}} = 0.75$, $x_{\text{crest,end}} = 0.8$,
    \Item{} Abyss begins at $x_{\text{abyss,start}} = 1$.
\end{Itemize}

We model transition zones between these regions using a smoothstep operator \cite{Perlin2002}:
\begin{align*}
    \smooth(x) = 3x^2 - 2x^3.
\end{align*}

For conciseness, denote the interpolating function as:
\begin{align*}
    S(a, b, x_0, x_1, x) = a + (b-a) \smooth\left(\frac{x - x_0}{x_1 - x_0}\right).
\end{align*}

The complete coral height field, as displayed in \cref{fig:coral-island-reef-function}, is built as a piecewise function:
\begin{align}
    \eqlabel{eq:coral-island-full-reef-function}{Analytic reef function}
    \heightCoral(x) &= \sum_{r \in \text{subregions}}{
    \begin{dcases}
        h_r & \text{if } x_{\text{r,start}} \leq x \leq x_{\text{r,end}} \\
        0 & \text{otherwise}
    \end{dcases}
    } \nonumber \\
    &+
    \sum_{t \in \text{transitions}} {
    \begin{dcases}
        S(h_{t}, h_{t+1}, x_{\text{t,end}}, x_{\text{t+1,start}}, x) & \text{if } x_{\text{t,end}} < x < x_{\text{t+1,start}} \\
        0 & \text{otherwise}
    \end{dcases}
    }.
\end{align}

\subsubsection{Blending height fields}
\label{subsubsec:height-functions-blending}

\begin{figure}
    \includegraphics[width=.45\linewidth]{smoothmax_min-max_k-5_large.png}
    \includegraphics[width=.45\linewidth]{smoothmax_min-max_k-50_large.png} \\
    \includegraphics[width=.45\linewidth]{smoothmax_min-max_k-5_short.png}
    \includegraphics[width=.45\linewidth]{smoothmax_min-max_k-50_short.png}
    \setfitratios{1, 1}
    \autofitcaptions[]{$k=5$, $k=50$}
    \caption[Definition of $\smoothmax$ as the average of $\smoothmax^+$ and $\smoothmax^-$]{The $\smoothmax^+$ operator (orange curve) is a function that is used to overestimate the maximum value of two input functions (dotted curves), especially when the difference between the two functions is small, while the $\smoothmax^-$ function (blue curve) tends to underestimate the $\max$ operator. Taking $\smoothmax$ (green curve) as the average of $\smoothmax^-$ and $\smoothmax^+$ creates a smooth function closely approximating the $\max$ operator, even with low values of sharpness $k$ (Left: $k=5$, right: $k=50$). Bottom graphs show a zoom on the domain $x \in [0.5, 1.3]$.}
    \label{fig:coral-island-blend-function-island-with-upper}
\end{figure}

\begin{figure}
    \autofitgraphics[]{smoothmax_all-k_large.png, smoothmax_all-k_short.png}
    \caption[Effect of $k$ on the $\smoothmax$ operator]{Blending two functions $f: \R \to \R$ (black, dotted) and $g: \R \to \R$ (blue, dotted) with the $\max$ operator (red curve), causing discontinuities when $f(x)=g(x)$, and with the $\smoothmax$ operator (green curves) with various sharpness values $k$, resolving the issue at the cost of slight overestimations with low values of $k$. Right presents a zoom on the domain $x \in [0.5, 1.3]$. }
    \label{fig:coral-island-blend-function-island}
\end{figure}

The final step is to blend the subsided height field $\heightSubsid(\p)$ with the coral feature height field $\heightCoral(\p)$ to produce the final terrain. The goal is to ensure that coral features remain near the water surface while allowing the rest of the island to subside.

To achieve this, we use a smooth maximum function defined as the mean of an overestimation and an underestimation of the $\max$ operator: 
\begin{align}
    \eqlabel{eq:coral-island-smoothmax-function}{Smoothmax function}
    \smoothmax(a, b) = \begin{dcases}
        a + \frac{1}{k} & \text{if } a = b, \\
        a + \frac{b - a}{1 - e^{-k(b-a)}} & \text{otherwise,}
    \end{dcases}
\end{align}
which smoothly blends the two height fields, ensuring that the coral regions dominates where coral growth is present, while the subsided island terrain dominates in other regions. This blending method ensures that the transition between the coral and subsided regions is smooth and visually consistent.

% We define our smooth maximum function $\smoothmax(a, b) \in \R$ for $(a, b) \in \R^2$ as the mean of two functions, $\smoothmax^-$ and $\smoothmax^+$ (shown in \cref{fig:coral-island-blend-function-island-with-upper}), adapted from Inigo Quilez's smooth min function \cite{Quilez2013}, that respectively underestimate and overestimate the function $\max$:

Here, $a = \heightSubsid(\p)$ is the height from the subsided island, $b = \heightCoral(\p)$ is the height from the coral reef feature, and $k$ controls the smoothness of the transition (fixed arbitrarly here at $k=5$ for heights ranging between 0 and 1). Higher values of $k$ bring the $\smoothmax$ function closer to the $\max$ function (\cref{fig:coral-island-blend-function-island}).

This smooth max function guarantees continuity by preventing abrupt height and slope differences between coral regions and the subsided terrain, creating a smooth, gradual transition that mimics the natural blending of coral reefs with deeper areas. The coral feature height field takes precedence where coral can grow, typically in shallow regions. In deeper regions, such as the abyss, the subsided height field naturally dominates, ensuring that the final terrain accurately reflects both subsidence and coral growth processes (\cref{fig:coral-island-volcano-example}).


\begin{algorithm}[H]
    \caption{Subsidence, coral growth, and smooth blending.}
    \label{alg:subsidence-coral-blend}
    \DontPrintSemicolon
    \KwIn{Deformed height field $\Tilde{h}$, subsidence rate $\subsidRate \in [0,1]$, coral growth function $\heightCoral$ (piecewise with transitions via $\smooth$)}
    \KwOut{Final height field $\hat{h}$}
    \BlankLine

    \tcp{(1) Subsidence}
    \ForEach{$\p$}{ $\heightSubsid(\p) \gets (1-\subsidRate)\cdot \Tilde{h}(\p)$ }

    \tcp{(2) Coral growth evaluation}
    \ForEach{$\p$}{
        \tcp{Use same $\distRegions_\p$ as in \cref{alg:sketch-to-height-wind}}
        $x \gets \distRegions_\p - i_{\text{reef}}$ \tcp*{$i_{\text{reef}} =$ reef region ID}
        $c(\p) \gets \heightCoral(x)$
    }

    \tcp{(3) Smooth maximum blend}
    \ForEach{$\p$}{ $\hat{h}(\p) \gets \smoothmax\left(\heightSubsid(\p),\, c(\p)\right)$ }

    \Return $\hat{h}$
\end{algorithm}

\subsubsection{Output}
\label{sec:coral-island-procedural-output}

\begin{figure}
    \autofitgraphics[]{volcano-input.png, volcano-render.png}
    \autofitgraphics[]{volcano-heightmap.png, volcano-corals.png, volcano-output.png}
    \autofitcaptions{Initial height field, Reef only, Blended height field}
    \caption[Volcano with coral reef generated using our curve-based method]{Volcano with single vent. Bottom left: the initial height field is computed directly from the user input. Bottom centre: the reef height field is output from \cref{eq:coral-island-full-reef-function}. Bottom right: blended result with \cref{eq:coral-island-smoothmax-function}.}
    \label{fig:coral-island-volcano-example}
\end{figure}

The resulting terrain represents a plausible coral reef island, where the volcanic island has subsided, and coral reefs have grown upward to keep pace with the water level. The smooth blending between the subsided terrain and the coral features ensures a natural transition between regions (island, lagoon, coral reefs, and abysses).

One of the key strengths of this method is its flexibility, as the subsidence and coral reef growth processes are modelled independently, allowing for a wide range of configurations. Users can generate plausible island terrains with or without coral features, or apply the coral reef growth modelling to existing height fields from other sources. Note that each of the pseudo-code section proposed in \cref{alg:subsidence-coral-blend} can be replaced with surrogate methods if desired (e.g., realistic subsidence or coral reef simulations, or the use of another blending function).
